\author{OTS}
\documentclass[fontsize=16pt]{mybib}
\renewcommand{\DFA}{12mm}
\renewcommand{\DZA}{1.3}

\graphicspath{{../../../assets/images/}}
\setheaderlogo{mnr-green.png}

% Absatzformatierung
\setlength{\parindent}{0pt}
\setlength{\parskip}{0.6em}
\renewcommand{\arraystretch}{0.8}

\title{Einleitung zu 1. Petrus}
\author{Lothar Schmid}
\date{}

\begin{document}
\mytitlepage
\tableofcontents
\newpage

\section{Teil 0: Einleitung}
Ich möchte mit euch zusammen die Briefe das Apostel Petrus durchgehen. Von Petrus haben es zwei Briefe in den Kanon der Bibel geschafft. Es gibt noch andere Schriften, die  Petrus zugeschrieben werden, aber diese wurden nie als kanonisch anerkannt. Wieso habe ich diese beiden Briefen gewählt?

Erstens, weil mich Petrus als Person fasziniert. Er ist einer wie du und ich, der gerade zur rechten Zeit am richtigen Ort war, um von Gott auserwählt zu werden, um an seinem Heilsplan direkt mitarbeiten zu dürfen.

Zweitens, die Briefe des Petrus passen sich nahtlos dem Evangelium von Jesus Christus und dem Evangelium des Paulus ein. Es ist quasi eine Zusammenfassung von den Paulus Briefen.

Drittens, diese Briefe werden für uns immer aktueller. Der wahre von der Schrift bezeugte Glaube, wird immer mehr in Frage gestellt, oder an weltliche Bedürfnisse angepasst.

Nun wollen wir zusammen diese Kostbarkeiten durchgehen. Wir fangen mit dem 1. Petrusbrief an und ich will klassisch vorgehen, wie ich es in der Schule gelernt habe. Zuerst verschaffen wir uns eine Übersicht des Briefes.\\
Folgende Punkte habe ich mir für Heute notiert:
\begin{enumerate}
    \item Der Verfasser
    \item Ort und Zeit der Abfassung
    \item Petrus der Apostel
    \item Inhalt des Briefes
\end{enumerate}

\section{Teil 1: Autor}
\biblezit{IPetr}{1:1} Der Verfasser des Briefes wird im ersten Vers erwähnt. Petrus schreibt von sich als Petrus, Apostel Jesu Christi. Nach dem Namen des Verfassers hat man diesen Brief den 1. Petrusbrief genannt.

Es gab eigentlich in der alten Kirchengeschichte keine Disskusionen ob, dieser Brief in den Kanon gehört oder nicht. Er ist schon sehr früh in den Sammlungen des Neuen Testaments aufgetaucht. \betonung[p]{(Ab Wann?)}

Erst in neueren Zeit wurde heftiger darüber gestritten, ob dieser Brief nun von Petrus geschrieben wurde oder nicht. Ein Kriterium war oder ist, dass dieser Brief in einem viel zu klassischen Griechisch geschrieben wurde, um von einem ungebildeten Fischer verfasst zu sein.\\
Wir lesen aber im Kapitel 5 Vers 12, dass ihm beim Verfassen des Briefes Silvanus, auch bekannt als Silas, zur Seite gestanden hat.
\biblezit{IPetr}{5:12-13}
Silas wird als Führender einer Gemeinde
\biblezit{Apg}{15:22}
und als Prophet bezeichnet
\biblezit{Apg}{15:32}
In der Anfangszeit der Gemeinde arbeiteten Apostel und Propheten eng miteinander. \biblezit{Eph}{2:20}

Mal ehrlich, wenn ihr dieses Predigtmanuskript vor der Korrektur meiner Frau zum Lesen bekommen hättet, wäre, jedem Deutschverständigen die Haare zu Berge gestanden. Meine Frau Janina hat das Manuskript korrigiert, aber diese Korrekturen haben nichts an dem Inhalt der Predigt geändert. 

So konnte Silas und oder vielleicht auch Markus, die ja beide zu der Zeit, als Petrus diesen Brief geschrieben hat, bei ihm waren, diesen Brief korrigieren, ohne etwas von dem Sinn des Inhalts zu ändern. \biblezit{IPetr}{5:12-13}

\section{Teil 2: Ort und Zeit der Abfassung}
Die Abfassungszeit wird auf das Jahr 64 \nChr{} oder 65 \nChr{} datiert. Es war in der Zeit als Nero Kaiser in Rom war. Nero ist ja einer der bekanntesten Kaiser der Römer und soll Rom angezündet haben. Seine Gegner gaben ihm die Schuld an dem Brand in Rom und Nero gab den Christen die Schuld und liess sie verfolgen und grausam töten.\\
Am Anfang war die Verfolgung eher auf Rom beschränkt und hat sich erst später über das Land ausgeweitet.

Mit der Abfassungszeit sind sich die Experten einig. Mit dem Abfassungsort wird es spannend. Aber lassen wir zuerst Petrus zu Wort kommen, wo er diesen Brief geschrieben hat.\\
Schlagen wir dazu \biblezit{IPetr}{5:13} auf. Da lesen wir, dass Petrus diesen Brief in Babylon geschrieben hat. Ist doch eigentlich alles klar so weit und wir können zum nächsten Punkt übergehen.

Aber nein, so klar ist das nicht. Es gibt zwei Lager, das eine Lager sagt, dieses Babylon sei ein Deckname für Rom und die andere Gruppe sagt, dieses Babylon sei das wirkliche Babylon in Mesopotamien im heutigen Irak.

Die Argumente der ersten Gruppe sind, dass Petrus diesen Namen verwendet hat, um die Christen in der Verfolgung zu schützen. Da Rom ein Sündenpfuhl war, konnten die Empfänger des Briefes leicht erkennen, dass dieser aus Rom kam. Der Deckname Babylon für Rom wird auch in der Offenbarungbenutzt. Diese Ansicht wird von den meisten Auslegern vertreten.

Die zweite Gruppe ist für das ursprüngliche Babylon. Ihre Argumentation: in Babylon gab es zu dieser Zeit eine grössere jüdische Gemeinschaft, bei der Petrus hätte zu Besuch sein können. Es gibt aber keine Hinweise, dass Petrus wirklich in Babylon war, aber auch nicht wirklich stichfeste Beweise, dass Petrus in Rom war. Ein bekannter Ausleger ist \zB{} Roger Liebi, der diese Ansicht vertritt.

Es kann sich jetzt jeder selbst Gedanken machen, welcher Gruppe er sich anschliessen möchte. Egal wo der Brief geschrieben wurde, dem Inhalt tut das absolut keinen Abbruch.

Die Überlieferung sagt, dass Petrus seinen Märtyrertod in Rom erlitten hat und darum wird er wohl auch irgenwann in Rom gewesen sein müssen.


\section{Teil 3: Der Apostel Petrus}
Um diesen Brief besser einordnen zu können, sollten wir uns auch mit diesem Apostel Petrus ein bisschen näher beschäftigen. Ich weiss, es wurde schon viel über diesen Petrus geredet und geschrieben, aber trotzdem möchte ich, bevor wir durch diesen Brief gehen, Petrus nochmals näher betrachten.

\subsection{Apostel}
Petrus nennt sich in der Anrede ein Apostel Jesu Christu \biblezit{IPetr}{1:1}. Apostel heisst, ein Abgesandter oder Botschafter. Auch wenn es Organisationen gibt, die ihre Leiter Apostel nennen, sind diese nach Johannes, dem letzten der 12 Apostel, ausgestorben. Dazu gibt es einen stichhaltigen Beweis in der Apostelgeschichte, der mich überzeugt, dass es keine weiteren neuen Apostel gegeben hat: 

Der Verräter Judas hat noch vor Pfingsten Selbstmord begangen. Der Heilige Geist war noch nicht ausgeschüttet und die restlichen elf Apostel waren der Meinung, dieser Judas muss ersetzt werden. Dieser Ersatz musste spezielle Kriterien erfüllen, \biblezit{Apg}{1:21-22}. Aus der Gruppe der Kandiaten wurden zwei ausgewählt, Joseph, der Barsabas und Matthias. Diese wurden nach mosaischem Gesetz mit dem Losverfahren ausgewählt. \biblezit{IIIMos}{16:8-10}.

Paulus, der Sonderapostel, wurde nicht von den Aposteln ausgewählt und musste nicht die gleichen Kriterien erfüllen. Paulus wurde von dem erhöhten Jesus Christus höchst persönlich auserwählt. Sowie die anderen 11 Apostel persönlich vom Menschen Jesus ausgewählt wurden. Gott ist halt immer noch der, der die Regeln macht.

Als Jakobus, der Bruder von Johannes, als erster Apostel den Märtyrertod erlitt \biblezit{Apg}{12:1-2}, wurde dieser nicht mehr ersetzt. Wenn die Linie der Apostel hätte weiter geführt werden sollen, dann hätte man ja auch für Jakobus wieder einen Nachfolger wählen müssen. Es wird aber in der Bibel nie mehr von einer neuen Apostelwahl geredet. Auch die frühen Kirchenväter erwähnten keinen neuen Apostel mehr.

\subsection{Petrus}
Zurück zu Petrus. Wie ihr sicher alle wisst, ist Petrus nicht sein ursprünglicher Name. Der Name Petrus wurde ihm von Jesus gegeben, und ist griechisch und bedeutet Stein oder eher kleiner Felsen. \biblezit{Matt}{16:18} Sein richtiger Name war Simeon.

Petrus stammte aus Bethsaida \biblezit{Joh}{1:44} und ist später nach Kapernaum umgezogen. Petrus war verheiratet \biblezit{Mk}{1:30} und auch Paulus erwähnt in seinem Brief an die Korinther \biblezit{IKor}{9:5}, dass Petrus seine Frau auf seinen Missionsreisen mitgenommen hat.

Petrus wird in den Evangelien und auch später in der Apostelgeschichte als ein Führer der Zwölf dargestellt. Bei den Auflistungen der Apostel wird er immer als erster erwähnt. Bei allen Schlüsselereignissen mit Jesus war Petrus immer mit dabei. \betonung[p]{Beispiele}.
In der Apostelgeschichte lesen wir, wie er auch danach wieder, von Anfang an die Führung, übernahm. \betonung[b]{Beispiele}.

In \biblezit{Matt}{16:18} lesen wir, dass Jesus, Petrus den Schlüssel zu dem Reich der Himmel gegeben hat. Jesus ist der Fels, auf dem die Gemeinde gegründet ist und Petrus ist quasi der erste Stein. Petrus ist dann der, der die ersten lebendigen Steine zu dieser Gemeinde hinzugefügt hat. Quasi, der Türöffner war. 
\begin{enumerate}
    \item Nach der Pfingstrede von Petrus \biblezit{Apg}{2:14-36} wurden auf einem Schlag 3000 Juden dieser neuen entstehenden Gemeinde Jesu hinzugefügt. Kann mir vorstellen, dass viele von denen von Taten Jesu gehöhrt und sogar miterlebt haben. 
    \item Auch wenn Paulus als der Heidenapostel bekannt ist, war nicht Paulus der die ersten Heiden zu Jesus führte, sondern Petrus den Hauptmann Kornelius. \biblezit{Apg}{10:}
\end{enumerate}
Petrus nimmt auch in seinem Brief auf den Vergleich mit den lebendigen Steinen Bezug \biblezit{IPetr}{2:6-8} 

Wenn wir das Leben von Petrus ab seiner Berufung von Jesus Christus ansehen, sehen wir, dass sich Petrus von Anfang an, sich für diesen Jesus aus Nazareth und seiner Lehre einsetzte. Er hat sein Geschäft verlassen und ist diesem Wanderprediger gefolgt, ohne zu Wissen, wohin ihn dieser Weg führen wird. Er kannt ja diesen Jesus nur vom hörensagen. Von den Predigten von Johannes des Täufers.\\
Drei Jahre war er mit Jesus und seinen Freunden auf Wanderschaft. Wer schon mal in Israel war, weiss wie hügelig die Gegend dort ist. Nur schon von Jerusalem bis Jericho sind es ca. 1100\~hm auf 25 km. Jesus ist in seiner kurzen Dienstzeit eine ordentliche Strecke gelaufen. Experten sagen bis zu 5000 km. Was mindestens 40 km pro Woche ist und oft von vielen Menschen begleitet.\\
Petrus war die ganze Zeit mit Jesus unterwegs. Bei Tag und Nacht, bei Hitze und Kälte, Hunger und Durst, aber auch bei schönen Sonnenuntergängen, Lachen und Lagerfeuerromatik und mit vielen Disskusionen und Gesprächen. \\
Wie das wohl war? So mit Jesus unterwegs zu sein? Johannes sagt von Jesus in seinem Brief \biblezit{IJoh}{3:4-5} die Sünde ist die Gesetzlosigkeit in Vers 4 und in Vers 5, in Jesus ist keine Sünde.\\
Könnt ihr euch das vorstellen? Mit jemanden unterwegs zu sein, der die Sünde nicht kannte wie Paulus in \biblezit{2Kor}{5:21} ausdrückt? Er war nie mürrisch, nie verärgert, nie aufbrausend, hat am Abend am Lagerfeuer nie über diese blöden Pharisäer gelästert. Er war immer zuvorkommend, geduldig und bescheiden. Und das in jeder Lebenslage und bei jeder Witterung. Es heisst ja, wir sollten immer, wenn wir etwas machen fragen, \enquote{was wohl Jesus tun würde?} Bei mir ist es eher so, dass ich nach einem Ausraster erst frage, was hätte Jesus getan. Dann ist es aber schon zu spät.

Als die Zeit Jesu hier auf Erden vollendet war, kam der Schmerz, die Verfolgung, das Leidens seines Lehrers und Weggefährten dazu. Wieso gerade mein Freund Jesus?  Der, der keine Sünde kannte? Der, der keiner Fliege etwas zu leidetat \biblezit{IPetr}{3:22-23}? Der, der alle Menschen geheilt hat, die zu ihm kamen \biblezit{Mk}{1:34}? Fragen auf die er damals noch keine Antwort hatte.

Obwohl Petrus diese enge Beziehung zu Jesus hatte, der mit Jesus im wahrsten Sinn des Wortes, durch dick und dünn gegangen ist, spricht er Jesus in seinem Brief nie als Kumpel oder Freund an. Sondern immer höchst respektvoll mit Jesus Christus, also mit dem vollständigen Titel Jesus, der Gesalbte an. Mit diesem Titel unterstreicht Petrus in seinen Briefen die Göttlichkeit von Jesus.

Wie ist es mit uns? Klar können wir Jesus als unseren Helfer und Freund ansehen. Wir dürfen und sollen ihm auch alle unsere Sorgen und Nöte darbringen, aber wir sollten nie vergessen, dass Jesus der Gesalbte ist. Der souveräne Schöpfer von Himmel und Erde den Johannes in seinem Evangelium \enquote{Logos} nennt \biblezit{Joh}{1:1-3}.\\
Darum sollten auch wir, unseren Herrn Jesus Christus, der für dich und mich am Kreuz gestorben ist, mit Respekt gegenüber treten. Auch Johannes, der auch mit Jesus 3 Jahre unterwegs war und sich sogar als Lieblingsjünger bezeichnet, begnet Jesus in seinen Schreiben immer äusserst respektvoll. Aber zurück zu Petrus.

Petrus war eine Führernatur. Es gibt solche Leute, die haben das im Blut. Kennt ihr die? Man soll in einem Kurs eine Gruppenarbeit erledigen oder man hilft einem Freund beim Umzug (zügeln). Zu Beginn sitzen alle herum und schauen sich gegenseitig an und wissen nicht, wo anfangen, bis dann einer aufsteht und die Führung übernimmt und anfängt zu leiten und Befehle zu erteilen. Diese Personen lassen sich aber oft auch nicht gerne korrigieren.

Petrus und seine Familie hatten einen Fischereibetrieb am See Genesareth \biblezit{Matt}{4:18-19}. Ich denke, Petrus hat diesen Betrieb nie aufgegeben, weil er nach der Beerdigung von Jesus wieder mit ein paar Kollegen auf den See fischen ging \biblezit{Joh}{21:2-3}. Hätte er alles verkauft, wäre er nicht so schnell an ein Schiff gekommen, um fischen gehen zu können. Petrus hatte sicher auch Angestellte, die während seiner Abwesenheit den Betrieb weiter am Laufen hielten. Seine Eltern und Verwandtschaft mussten ja auch von irgendwas leben. Sind jetzt nur so gedanken von mir. In der Bibel lesen wir ja nichts davon.

Petrus war mutig und sicher auch abenteuerlustig, aber er war ganz sicher nicht naiv und hat alles hingeschmissen und seine Liebsten im Stich gelassen. Trotz seiner Abwesenheit wurden seine Familie und seine Angestellten weiterhin versorgt. Es zeigt auch, das Petrus ehrlich im Geschäft war und so dieses einfach vertrauensvoll abgeben konnte. 
Eigenschaften, die wir auch in seiner Nachfolge mit Jesus erkennen können, sowie auch in diesem Brief, den wir hier gemeinsam anschauen werden.

Petrus war ein Mensch wie du und ich, mit Ecken und Kanten, aber auch mit Eigenschaften, die zur Nachfolge wichtig sind. Mut, Ehrlichkeit und Loyalität. Eigenschaften die auch in der Ehe, Familie, am Arbeitsplatz und in der Gemeinde wichtig sind. Petrus war mutig genug um seinen Betrieb zu verlassen, um Jesus nachzufolgen, Petrus war ehrlich, \zB{} als Jesus seine Leiden ankündigte, nimmt Petrus ihn beiseite und fing an ihn zu wehren \biblezit{Mk}{8:32}. Petrus wahr ehrlich zu Jesus und hat ihm seine persönliche Meinung gesagt.\\
Er war loyal. Petrus war bis zuletzt bei Jesus. Als Jesus vor Pilatus geführt wurde, war Petrus noch bei ihm. Nach dem Hahnenschrei, als ihm bewusst wurde, dass er Jesus verraten hatte, hat es ihn bitterlich bereut \biblezit{Matt}{26:75}. Nein, Petrus war nicht vollkommen, er hat nicht alles richtig gemacht. Er hat aber gelernt und hat sich weiter entwickelt und das sehen wir in seinen Briefen.

Wenn du ein Nachfolger von Jesus Christus bist, brauchst auch du diese Eigenschaften. Eigenschaften, die für das tägliche Leben förderlich sind und die wir für den Bau der Gemeinde Jesu einsetzen können und sollen. Es gelingt nicht immer so wie es sein sollte, aber wie bei Petrus wird auch dir Jesus Christus helfen diese Eigenschaften zu lernen und dich darin zu festigen.

\begin{enumerate}
    \item [Mutig] Mutig sein heisst, dein ganzes Vetrauen auf Gott zu setzen, in der Gewissheit, dass er alles gut machen wird. Du sollst ihn nicht versuchen, aber du kannst dich mutig für die Belange Gottes einsetzen. Zum Beispiel bei Gesprächen ihn Bezeugen und von und über ihn reden. 
    \item [Ehrlichkeit]Sei ehrlich gegenüber Gott. Gott will, dass du mit allem zu ihm kommst. Es nützt ja auch nichts etwas vor ihm zu verbergen, aber ich \zB ertappe mich immer wieder, dass ich etwas verschweigen will, oder denke ich sage es ihm später. Gewöhnlich geht mir dann ein Lichtlein auf, wie blöd, Gott weiss ja alles. Ehrlich sein, auch gegenüber seinen Mitmenschen. In der Gemeinde, Familie, am Arbeitsplatz oder generell unter Menschen, seid einfach ehrlich.
    \item [Loyalität] Loyalität gegenüber Jesus Christus. Zu Beginn meiner Laufbahn mit Jesus tat ich mich schwer. Heute freue ich mich, wenn ich angesprochen werde. Ich bin nicht der Evangelist der auf Menschen zu geht und sie einfach so anspricht und über Jesus redet, aber wenn mich jemand auf meinen Glauben anspricht, rede ich gerne über Jesus und nutze die Gelegenheit \biblezit{Matt}{10:32-33}.
\end{enumerate}

Wie wir am Werdegang von Petrus sehen, wurden diese Eigenschaften von Jesus und danach vom Heiligen Geist in ihm gefördert. So kannst auch du gewiss sein, dass der Heilige Geist dir in dieser hinsicht seine Hilfe gibt und dich unterstützt.

Petrus ging mit seiner Loyalität sogar so weit, dass er für seinen Herrn Jesus Christus als Märtyrer in Rom gestorben ist. Laut Überlieferung wurde er von Nero zum Tod am Kreuz verurteilt und weil Petrus sich zu gering ansah, gleich zu sterben, wie sein Herr Jesus, war sein Wunsch, kopfüber gekreuzigt zu werden. In der Acta Petri oder Petrusakte (180-200\nChr), wird die Geschichte so erzählt:
\begin{quote}
    Petrus gerät in Rom mit Simon Magus einem Magier in den Streit. Unter dem Verfolgungsdruck will Petrus Rom verlassen. Auf dem Weg begenet ihm Jesus. Petrus fragt Jesus: \enquote{Quo vadis, Domine?} (Wohin gehst du, Herr?).\\
    Christus antwortete ihm, er gehe nach Rom, um sich erneut kreuzigen zu lassen. Petrus versteht diese Antwort als Ruf, dass er Umkehren soll. Er kehrt also nach Rom zurück und wird kopfüber gekreuzigt.
\end{quote}
Nur eine Geschicht und wenn ich ehrlich sein will, nicht unbedingt glaubwürdig. Es gibt keine handfesten Beweise, dass Petrus je in Rom war. Die Überlieferung alter Schriften sagt das. Wenn er der Gründer der katolischen Kirche sein soll, dann muss er ja in Rom gewesen sein.

Petrus war aber sicher nie, der Gründer der Gemeinde in Rom und er war auch nie Bischof in Rom und ganz sicher nicht Papst. Er war auch nicht zusammen mit Paulus im Rom, weil Paulus schreibt in \biblezit{IITim}{4:11} das nur Lukas bei ihm in Rom ist. Ausserdem taucht er in keiner Grussliste in den Briefen von Paulus auf, die er an Rom oder in Rom geschrieben hat.

Vielleicht hat Petrus diesen Brief doch in Babylon geschrieben. \betonung[b, p]{Augenzwinkern}


\section{Teil 4: Inhalt des Briefes}
Ich habe die Kapitel mit den folgenden Überschriften versehen:

\begin{enumerate}
    \item Leben als Christ in der Hoffnung
    \item Leben als Christ in der Gnade
    \item Leben als Christ in der Welt
    \item Leben als Christ unter Christen
    \item Leben als Christ in der Gemeinde
\end{enumerate}

Petrus schrieb diesen Brief an Christen in der Region von Kleinasien in der heutigen Türkei. \biblezit{IPetr}{1:1}. An sogenannte Juden-Christen in der Diaspora, was soviel heisst wie Zerstreuung.

Dieser Brief soll die Christen in dieser Region ermutigen, in ihren Leiden auszuharren. Bis zu 15--mal spricht er, mit 8 verschiedenen griechischen Begriffen, ihr Leiden an. Sowohl Christen als auch Juden wurden zu dieser Zeit verfolgt. So hat Kaiser Claudius von Rom um 53 \nChr{} alle Juden aus Rom verjagt und umgesiedelt. Dazu gehörten auch Juden, die Christen geworden waren.

Da die Abfassungzeit des Briefes um 64 \nChr{} und der Brand in Rom auch erst um 64 \nChr{} stattfand, wird wohl nicht dieser Brand Schuld an der Verfolgung dieser Christen in diesen von Rom entfernten Regionen gewesen sein.

Wie dem auch sei, sicher war, dass Petrus sie mit diesem Brief ermutigen wollte. Dieser Brief war nicht nur eine Ermutigung für die Christen zur damaligen Zeit, sondern auch noch Jahrhunderte später. Bis in die heutige Zeit hinein ermutigt er uns und alle Christen die verfolgt werden.

Ich selbst habe noch keine Verfolgung erlebt, Spott oder Hohn ja, aber jetzt nicht wirklich schlimmes. Auf der Arbeit werde ich manchmal belächelt, aber wirklich Nachteile hatte ich bis jetzt nie.

Ich bin mir sicher, dass wir vor der Trübsal entrückt werden. Aber ich schliesse nicht aus, dass auch wir hier im Westen eine Zeit der Unterdrückung erleben werden. Schon jetzt wird ja der wirklich Bibeltreue Christ ein bisschen schräg angesehen.

Laut Open Doors ist die Christenverfolgung so gross wie noch nie seit es Christen gibt. Die Verfolgten, wie auch wir hier im noch wohlbehüteten Westen, gehören alle dem gleichen Leib Christi an. Wie Paulus in 1. Korintherbrief schreibt, dass, wenn ein Glied leidet, die anderen Glieder mitleiden. Mit unserem Wohlstand können wir durch Gebet und Gaben die verfolgten Christen unterstützen.\\
Petrus hat zur Unterstützung der leidenden Christen einen Brief geschrieben, der sie ermutigen soll, trotz allem auf ihren Erlöser Jesus Christus zu schauen. Petrus hatte sicher auch noch die Worte von Jesus im Kopf, wo Jesus ihm 3x sagte: \enquote{Hüte meine Schafe!} Mit diesem Brief wollte er diesem Auftrag gerecht werden.\\
Auch Open Doors bietet die Möglichkeit an, ermutigende Briefe an verfolgte Christen zu schicken. Man verfasst einen Brief und sendet ihn Open Doors. Diese verändern ihn dann so, dass er für den Empfänger keine Gefahr mehr ist und sendet versendet sie dann.

Beim Studieren dieses Briefes sind mir die vielen Paralellen zu den Paulusbriefen aufgefallen. Etwas, das mich sehr erfreut. Ich bin immer wieder erstaunt, wie alles in diesem ganzen Buch so harmonisch zusammenpasst, obwohl soviele Schreiber über einen so langen Zeitraum daran gearbeitet haben. Wir dürfen nicht vergessen, Paulus und Petrus haben sich selten gesehen. Es gab noch kein Handy, um sich gegenseitig abzusprechen und auszutauschen. Sie wurden beide, wie auch alle andere Autoren, die an diesem Buch geschrieben haben, vom Heiligen Geist geführt.

Die Gnade Gottes ist für Petrus sehr wichtig. Die Gnade kommt in jedem Kapitel vor. \biblezit{IPetr}{1:2.10.13}, \biblezit{IPetr}{2:19-20}, \biblezit{IPetr}{3:7}, \biblezit{IPetr}{4:10}, \biblezit{IPetr}{5:5.10.12}.

So wie Petrus hier die Gnade immer wieder betont, können auch du und ich uns immer wieder auf diese Gnade berufen.
\begin{enumerate}
    \item Gnade gibt Hoffnung \biblezit{IPetr}{1:13}
    \item Gnade rettet \biblezit{Eph}{2:8-10}
    \item Gnade gibt Kraft \biblezit{IIKor}{12:1-10}
    \item Gnade führt zur Herrlichkeit \biblezit{Ps}{84:11}, \biblezit{Ps}{5:10}
\end{enumerate}

Das heisst, wenn du Jesus Christus von Herzen aufgenommen hast, kannst du mit der Hoffnung, die wir daraus erhalten, dein Leben meistern. Mit dieser Hoffnung werden plötzlich deine irdischen Probleme klein. Die Hoffnung ist die Freiheit, die du als Christ hier auf dieser Erde hast.

Diese Hoffnung wird aber betrübt, wenn du bewusst in Sünde lebst. Um die Hoffnung wiederzuerlangen, musst du neu Busse tun und dein Blick auf das Kreuz richten. Die Gemeinde kann helfen, dass diese Hoffnung wieder erlangt werden kann. Die Gemeinde ist dazu da, sich gegenseitig zu ermahnen, zu unterstützen und zu helfen \biblezit{Hebr}{10:23-25}.\\
Es gibt da ein richtig blödes Sprichwort: \enquote{Die Hoffnung stirbt zuletzt.}\\
Nicht aber für dich und für mich als Christ. Unsere Hoffnung ist zuerst gestorben, nähmlich am Kreuz. Lass diese auferstandene Hoffnung in dein Leben. Du bist in ihm gestorben und mit ihm wieder auferstanden zum ewigen Leben. \biblezit{Rom}{6:3-4}.

\begin{block}[Ende]
    Ich habe euch jetzt einen kurzen Überlick die Person Petrus und über seinen ersten Brief gegeben.
    Petrus ist für mich persönlich eine besondere Person. Paulus ist für mich eher der unnahbare, der studierte, der intelektuelle, der spezielle Apostel. Paulus taucht einfach plötzlich auf. Petrus aber, ist quasi mit Jesus aufgewachsen, er gehört einfach dazu. Es ist spannend, wie diese zwei Briefe seine geistige Entwicklung so schön aufzeigt. Ich freue mich darauf mit euch zusammen diese Briefe durchzuarbeite und seine Kostbarkeiten zu ergründen. 
\end{block}

\end{document}